% SPDX-License-Identifier: GPL-3.0-or-later
% Copyright (C) 2026 Alessandro Greco
\section*{Глава \ref{cap:23-browser}: Tor Browser, Tails и Whonix}

\begin{soluzione}{es:23-identita}
Функция «Новая цепь для этого сайта» (New Circuit for This Site) меняет только цепь сайта активной вкладки, но не удаляет файлы cookie, историю и состояние других сайтов: сайт может по-прежнему узнавать пользователя по сохранённым данным. «Новая личность» (New Identity) закрывает вкладки, очищает состояние браузера и создаёт новые цепи для дальнейшей работы. Ни одна из команд не меняет входной узел (guard), поскольку частая замена повышает вероятность рано или поздно выбрать враждебный узел (\cref{sec:guard}); при этом входной узел не знает посещаемых сайтов.
\end{soluzione}

\begin{soluzione}{es:23-estensioni}
Дополнительное расширение может блокировать загрузку некоторых отслеживающих элементов, но Tor Browser уже изолирует их по сайтам и ограничивает устойчивые идентификаторы. Цена — потеря единообразия: сайты могут обнаружить само расширение или его списки фильтрации и выделить пользователя среди других пользователей Tor Browser. Кроме того, расширение добавляет программный код с доступом ко всем посещаемым страницам.
\end{soluzione}

\begin{soluzione}{es:23-isolamento}
Если две вкладки одного сайта используют одинаковые учётные данные SOCKS, tor направляет их потоки через одну цепь; третья вкладка с другими данными получает другую цепь. Команда \texttt{GETINFO stream-status} позволяет увидеть идентификатор цепи каждого потока.
\end{soluzione}

\begin{soluzione}{es:23-scelta}
(a) Достаточно Tor Browser, при необходимости с мостом. (b) Tails: загружается с USB-накопителя, не оставляет обычных следов на чужом компьютере и сохраняет только то, что помещено в зашифрованное постоянное хранилище; сохраняется риск подменённого оборудования. (c) Whonix: приложения Workstation проходят через tor по устройству системы, которая подходит для длительной работы без обязательного удаления следов при завершении сеанса.
\end{soluzione}

\begin{soluzione}{es:23-impronta}
Упражнение с открытым ответом. Примеры характеристик: размеры экрана и окна, часовой пояс, язык, доступные шрифты, результат отрисовки на холсте (canvas). В обычном браузере они могут выдавать от нескольких бит информации (например, часовой пояс) до многих (шрифты и canvas). Tor Browser округляет размеры окна с помощью letterboxing, сообщает UTC и единый язык, ограничивает шрифты и унифицирует либо блокирует наиболее раскрывающие интерфейсы. Остаются различия, связанные с операционной системой, уровнем безопасности и округлёнными размерами экрана.
\end{soluzione}

\section*{Глава \ref{cap:24-metrics}: измерение сети}

\begin{soluzione}{es:24-dieci}
Количество запросов к каталогам уменьшилось бы вдвое, а при постоянном делителе 10 вдвое упала бы и оценка числа пользователей. Графики показали бы почти одновременное снижение в разных странах по мере распространения новой версии tor. Постепенный глобальный спад, совпадающий с внедрением версии, свидетельствует об изменении поведения программы, а не числа пользователей; метод оценки нужно исправить, изменив делитель.
\end{soluzione}

\begin{soluzione}{es:24-frazione}
$20\,000\,000 / (0{,}6 \cdot 10) \approx 3\,333\,333$ пользователя. При доле отчётности 8\% оценка отбрасывается, поскольку она ниже порога 10\%.
\end{soluzione}

\begin{soluzione}{es:24-flag}
План: отберите ретрансляторы с обоими флагами и разделите суммарную пропускную способность этих узлов на общую пропускную способность сети. Для вероятности выбора выходного узла возьмите \texttt{probabilities(consensus, "exit")} из Лаборатории~7 и сложите вероятности ретрансляторов, чьи сети /16 входят в пятёрку наиболее часто встречающихся.
\end{soluzione}

\begin{soluzione}{es:24-censura}
Упражнение с открытым ответом. Обычно число прямых пользователей резко падает в день блокировки, а число пользователей мостов, в частности Snowflake или WebTunnel, может расти в следующие дни по мере перехода на подключаемые транспорты. Индикатор Metrics отмечает дату первоначального падения.
\end{soluzione}

\begin{soluzione}{es:24-privacy}
Без округления противник, знающий, что является единственным пользователем определённой страны на данном ретрансляторе (или точно знающий число других пользователей), может по статистике extra-info определить, появился ли ещё один пользователь этой страны: счётчик изменится на единицу. Округление скрывает такие небольшие изменения. Противник, контролирующий много ретрансляторов, вообще не нуждается в опубликованной статистике: он непосредственно видит подключения к своим узлам.
\end{soluzione}

\section*{Глава \ref{cap:25-legale}: правовые и этические аспекты}

\begin{soluzione}{es:25-condizioni}
(a) Выходной ретранслятор, самостоятельно отправляющий запросы на сайты, например для \enquote{смешивания} трафика, сам инициирует передачу. (b) Выходной ретранслятор, перенаправляющий соединение на адрес, отличный от запрошенного в ячейке \texttt{BEGIN}, выбирает получателя. (c) Выходной ретранслятор, вставляющий содержимое в страницы HTTP или фильтрующий их по смыслу, выбирает либо изменяет информацию. Именно за подобное поведение центры каталогов могут присвоить флаг \texttt{BadExit}.
\end{soluzione}

\begin{soluzione}{es:25-strumento}
Кухонный нож как инструмент законен, нападение с ним — нет. Чтение запрещённой цензурой газеты с помощью Tor может быть законным применением законного инструмента, но неправомерный доступ к компьютеру остаётся правонарушением и с Tor, и без него. Оператор ретранслятора предоставляет канал передачи, не выбирая содержимое: он ближе к перевозчику, чем к автору, если сам не участвует в незаконной деятельности. Конкретные правовые выводы зависят от страны и обстоятельств.
\end{soluzione}

\begin{soluzione}{es:25-segnalazione}
Следует объяснить, что адрес принадлежит выходному ретранслятору Tor, что проверяется публичными инструментами Tor Project; что оператор не знает и не записывает личность клиентов и не может их установить; что сайт при желании может отдельно обрабатывать подключения с публичных выходных адресов. На практике оператор может добавить в политику выхода запрет для адреса жалующегося сайта, если польза оправдывает ограничение.
\end{soluzione}

\begin{soluzione}{es:25-esperimento}
Предложение предполагает наблюдение за трафиком реальных пользователей без их согласия, сбор сведений, которые нельзя безопасно публиковать (посещённые домены), и лишь небольшую пользу. Оно нарушает несколько принципов безопасного исследования и в Италии может создавать проблемы уголовного права и защиты персональных данных. Альтернативы: использовать уже опубликованную агрегированную статистику Metrics, самостоятельно проверять доступность новостных сайтов через собственный клиент Tor либо предложить Совету по безопасности исследований (Research Safety Board) схему подсчёта с агрегацией и добавлением шума, как в статистике ретрансляторов.
\end{soluzione}

\begin{soluzione}{es:25-ricerca}
Исследовательское задание: ищите судебные решения в официальных базах (в Италии — решения Кассационного суда и других судов, в ЕС — документы Суда Европейского союза), а не только в сторонних пересказах. Для каждого решения укажите суд, дату, применённую норму и результат; различайте дела об операторах ретрансляторов и о пользователях Tor.
\end{soluzione}

\section*{Глава \ref{cap:26-attacchi}: корреляция, анализ отпечатков и Sybil-атаки}

\begin{soluzione}{es:26-posizione}
Корреляция трафика требует наблюдения за обоими концами соединения (ретрансляторами или сетями); защиты: входные узлы, ограничения маршрутов, взвешивание по измеренной пропускной способности. Определение посещаемого сайта по отпечатку трафика требует лишь наблюдения канала клиент—входной узел; защиты: ограниченное заполнение, исследования продолжаются. Обнаружение входного узла требует контроля промежуточного ретранслятора и возможности вызывать построение цепей; защиты: постоянные входные узлы и vanguards. Атаки Sybil используют множество ретрансляторов; защиты: измерение пропускной способности, семейства, ограничение /16 и наблюдение за сетью. Атаки маршрутизации требуют контроля автономной системы (AS) или точки обмена трафиком (IXP); внешняя защита — мониторинг BGP.
\end{soluzione}

\begin{soluzione}{es:26-base}
На 100\,000 посещений приходится в среднем около 0,95 правильного сигнала и $0{,}001 \cdot 99\,999 \approx 100$ ложных. Доля правильных сигналов составляет $0{,}95/(0{,}95+100) \approx 0{,}9\%$. Почти все предупреждения оказываются ложными.
\end{soluzione}

\begin{soluzione}{es:26-bin}
При двухсекундных интервалах за пять секунд остаётся только две или три точки на ряд, поэтому корреляция становится случайной. Если задать всем клиентам один ритм трафика (например, одинаковое зерно генератора случайных чисел), ряды становятся почти неразличимыми и сопоставление затрудняется. Именно выравнивание ритмов сильнее всего мешает корреляции: такова идея постоянной скорости передачи в PipeNet (\cref{sec:radici-bassa}).
\end{soluzione}

\begin{soluzione}{es:26-difesa}
При постоянном темпе передачи входные ряды перестают меняться. Корреляция Пирсона тогда математически не определена (нулевая дисперсия) либо определяется шумом сети: такая атака не работает. Цена — постоянная передача с максимальной выбранной скоростью, даже когда полезных данных нет. В лаборатории это десятки или сотни KB/s на клиента, намного больше полезного трафика.
\end{soluzione}

\begin{soluzione}{es:26-probabilita}
С постоянным входным узлом противник выигрывает, только если этот узел принадлежит ему (вероятность около 5\%); тогда примерно в 5\% цепей выходной узел тоже принадлежит противнику. После шести месяцев интенсивного использования вероятность хотя бы одной компрометации приближается к 5\% в целом, а при честном входном узле остаётся нулевой в этой упрощённой модели. Если выбирать входной узел заново для каждой цепи, риск одной цепи около $0{,}05 \cdot 0{,}05 = 0{,}0025$; после нескольких тысяч цепей хотя бы одна компрометация почти неизбежна. Это аргумент в пользу постоянных входных узлов; проверку моделированием позволяет класс \texttt{Adversary} из Лаборатории~7.
\end{soluzione}

\section*{Глава \ref{cap:27-casi-rete}: реальные атаки на сеть}

\begin{soluzione}{es:27-campo}
В направлении к клиенту каждый ретранслятор добавляет свой слой шифрования к содержимому ячейки (\cref{cap:10-cifratura-relay}). Входной узел видит содержимое, зашифрованное ключами промежуточного и собственного переходов, и не может прочитать данные, записанные HSDir. Команда самой ячейки, напротив, находится во внешнем заголовке, который каждый ретранслятор читает и пересылает без изменений.
\end{soluzione}

\begin{soluzione}{es:27-uno}
При трёх входных узлах вероятность, что хотя бы один из них враждебен, примерно втрое выше, и каждый такой узел видит часть трафика клиента. При одном входном узле риск выбора противника ниже; если узел честен, клиент защищён долгое время. Минус: единственный враждебный входной узел наблюдает весь трафик клиента, а медленный или недоступный узел снижает производительность.
\end{soluzione}

\begin{soluzione}{es:27-verso}
Лаборатория уже использует ячейки \texttt{DROP}, а не запросы на расширение цепи: счётчик в \file{command.c} проверяет внешнюю команду ячейки до расшифрования, поэтому не зависит от её содержимого. В обратном направлении программа записывает предупреждение и закрывает цепь с причиной \texttt{PROTOCOL}, независимо от того, принадлежит ли цепь клиенту или ретранслятору.
\end{soluzione}

\begin{soluzione}{es:27-segnali}
Признаки: ретрансляторы появились в консенсусе в одну минуту; в дескрипторах совпадают редкая версия tor и операционная система; адреса находятся в одной сети или у одного провайдера; похожи порты, имена узлов, отсутствуют контакты. Ложные совпадения возможны у честных крупных операторов, одновременно обновляющих узлы, распространённых хостинг-провайдеров и автоматических установок с одинаковой конфигурацией. Одного признака недостаточно: нужны разные свидетельства и, по возможности, связь с оператором.
\end{soluzione}

\begin{soluzione}{es:27-canali}
Промежуточный ретранслятор видит незашифрованные внешние данные: CircID (на каждом переходе меняется, поэтому через всю цепь не проходит); команду ячейки (\cella{RELAY} либо \cella{RELAY\_EARLY}): этот канал был использован в 2014 году, теперь он закрыт в направлении к клиенту и ограничен восемью ячейками в обратном направлении; ячейку \cella{DESTROY} и причину закрытия, время и число ячеек. Последние две величины всегда доступны для пассивного подтверждения и не имеют полной защиты. Содержимое зашифровано послойно, но в tor1 один злонамеренный ретранслятор может изменять его для распознавания другим (маркировка, \cref{sec:tor1-debolezze}); CGO (\cref{sec:cgo}) делает такие изменения причиной невосстановимого содержимого на следующих переходах.
\end{soluzione}
